% ===========================================================================
%  sviluppatore/07-marketplace/05-dati-ausiliari.tex
% ===========================================================================
\section{I dati ausiliari di un artefatto}\label{sez:dati-ausiliari}

Una strategia che legge gli utili trimestrali, senza quei file, gira lo
stesso: \file{core/node\_data.py} restituisce serie a zero invece di un
errore. Gira però su un segnale vuoto, e il backtest dell'autore non dice più
niente di ciò che succederà sul computer di chi compra. Per questo la scheda
prodotto elenca le categorie di dati che l'artefatto legge, con l'avviso che
senza quei dati non ha la garanzia di funzionare; la card della vetrina le
accenna.

\subsection{Da dove si leggono}

Dal \textbf{codice pubblicato}, non da una dichiarazione dell'autore:
\file{data/dati\_ausiliari.py} analizza i file con \pkg{ast}, senza
importarli (con un ripiego a espressioni regolari per il codice che non si
lascia analizzare), e cerca gli stessi segni che il sistema usa per
procurarsi i dati:

\begin{itemize}
  \item \cod{required\_external\_data=[...]} nei descrittori e la stessa chiave
    nei \cod{META} di overlay e weight mode;
  \item \cod{required\_mms\_features} diverso da \cod{None}: la categoria
    \cod{mms};
  \item le chiamate \cod{data\_series("earnings", ...)} e
    \cod{custom\_series(...)}, anche delle strategie a blocchi;
  \item gli import dei vecchi lettori (gli indizi di
    \file{data/external\_data.py});
  \item la cartella \cod{FOLDER} della porta sui dati (\file{data\_port.py}).
\end{itemize}

\noindent
Le serie calcolate dal motore sui prezzi (\cod{volume\_profile}) e le
cartelle di sistema non contano. Per un project i file sono quelli di tutti i
componenti. Ogni categoria porta etichetta e origine: \cod{sistema} (con la
fonte e le chiavi del \file{.env} lette da \file{data/fonti}), \cod{app} per
le MMS, \cod{personale} per quelle create dall'autore.

\subsection{Dove finiscono}

Con la pubblicazione in \cod{details.datiAusiliari} (con \cod{portaDati}):
un elenco vuoto vuol dire «controllato, nessuna», l'assenza del campo
«pubblicato prima del controllo». Per questi ultimi la scheda le ricava al
momento quando chi guarda vede il codice, e l'autore le salva aprendo la
propria scheda o \voce{My products} (\simbolo{web\_db.salva\_dati\_ausiliari},
che tocca solo artefatti propri). La vetrina le legge con
\cod{a.details->'datiAusiliari'}; finché il sito non conosce quella query,
\simbolo{web\_db.marketplace\_list} ripiega sulla precedente e la vetrina
resta piena, senza le categorie sulle card.

\noindent
La prova è \file{tests/prova\_dati\_ausiliari.py}: le strategie vere del
deposito, una strategia a blocchi compilata, la porta sui dati, la scheda
disegnata da Flask e la card eseguita con \pkg{node}.
